On laisse tomber une balle A, sans vitesse initiale, du haut d'un édifice de \qty{\v{H}}{m}. Au même instant, une balle B est lancée verticalement vers le haut depuis le sol, juste sous A, avec une vitesse de \qty{\v{v0}}{m/s}. On néglige la résistance de l'air.

\begin{enumerate}[label=\alph*)]
    \item À quel instant les deux balles se croisent-elles ?
    \item À quelle hauteur au-dessus du sol se croisent-elles ?
    \item Quelle est alors la vitesse de la balle B ? Monte-t-elle ou descend-elle ?
\end{enumerate}
